%
%
\documentclass[11pt]{article}
\usepackage[preprint]{acl}
%
\usepackage{times}
\usepackage{latexsym}
\usepackage[T1]{fontenc}
\usepackage[utf8]{inputenc}
\usepackage{microtype}
\usepackage{amsmath}
\usepackage{amssymb}
\usepackage{booktabs}
\usepackage{graphicx}
\usepackage{tikz}
\usetikzlibrary{positioning,arrows.meta,fit,backgrounds,calc}
\usepackage{enumitem}
\usepackage{xcolor}
\usepackage{listings}
\usepackage{upquote}

%
\DeclareUnicodeCharacter{0394}{\ensuremath{\Delta}}

%

%
\newcommand{\sysname}{our framework}

%


%

\title{Patching the Blind Spots: Reversible LLM Rule Overlay\\
       for Frozen Fraud Detectors}

%
  \author{%
     \textbf{Bo Qu\textsuperscript{\ensuremath{\dagger}}\thanks{Corresponding Author.}},
      \textbf{Zhurong Wang\textsuperscript{\ensuremath{\dagger}}},
    \textbf{Frank Zahradnik\textsuperscript{\ensuremath{\dagger}}},
     \textbf{Ming Ouyang\textsuperscript{\ensuremath{\dagger}}}
    %
    \\
     \textsuperscript{\ensuremath{\dagger}}eBay,
     \\
\texttt{{boqu@ebay.com}}
  }

\begin{document}
\maketitle
\begin{abstract}
In industrial tabular fraud classification, gradient-boosted trees remain state-of-the-art, but even
a well-trained model has blind spots, and retraining the deployed model is not always safe. We
propose a framework that patches such a model without modifying it: LLMs propose rules
anchored on features the model underweights by SHAP, and the resulting noisy pool is screened for
empirical support and calibrated into an additive, reversible
log-lift overlay that runs at the production cut-off. Single-LLM benefit is high-variance and inverts between datasets, so we pool three
diverse backbones by equal-weight logit averaging---trading peak gain for a bounded worst case. On a
real-world e-commerce fraud detector trained on
selection-biased data, continued training regresses even when carefully tuned, whereas the overlay
gains on average $+2.26\%$ recall at equal precision ($\Delta$R@P) and $+1.60\%$ precision at equal
recall ($\Delta$P@R) over production cut-offs spanning 30--57\% precision and 23--63\% recall. On a
public benchmark of seven designed bias-shift variants it raises $\Delta$P@R on every variant and
keeps $\Delta$R@P near neutral. Each amendment is a readable predicate naming the fraud pattern it
covers.
\end{abstract}

\section{Introduction}
\label{sec:intro}

Gradient-boosted trees dominate tabular fraud classification
\citep{ke2017lightgbm,chen2016xgboost}, trained once on historical labelled transactions, then
degrade under the temporal distribution shift ubiquitous in financial data
\citep{quinonero2009shift,gama2014drift,dalpozzolo2018fraud}: new fraud patterns emerge, the
population shifts, and discrimination quietly erodes with no visible retraining signal.

In production, such a detector also has \emph{loopholes}: fraud patterns it scores too low, which
persist after deployment. Closing them is constrained in two ways in our setting. The model's scores
and cut-offs drive downstream decisions, so changing the model means a redeployment. Its training
data are selection-biased: a rule engine and a screening model filter transactions before they are
recorded, so training sees \emph{approved} traffic only, while the model serves the full population
(a small unbiased cohort is kept for evaluation; Section~\ref{sec:setup}). Patterns that rarely pass
the filters are under-represented, and the usual fix---continued training on a fresh window---fits
the same bias: in our experiments it regresses on the unbiased cohort (Section~\ref{sec:results}).
The practical question is how to find a frozen model's loopholes and improve how it generalizes,
without modifying it.

\paragraph{Key idea: LLMs reason about the frozen GBM's blind spots and patch them proactively.}
Instead of updating the model, we let large language models reason over its diagnostics (SHAP
attributions and a digest of its trees) and propose human-readable rules for the fraud it misses; the
adaptation window is used only to screen and calibrate them, so no model parameter changes.
Generation is steered \emph{through the model's weak dimensions}: each rule's primary anchor must be a
feature the model underweights by SHAP. The rules are calibrated into a band-conditioned log-lift
overlay on the \textbf{frozen} base---additive in logit space, hence reversible---and mapped back to
the production score scale, so the overlay runs at the existing cut-off. Every amendment is a
human-auditable predicate, in line with calls for interpretable high-stakes models
\citep{rudin2019interpretable}.

Prior work uses LLMs on tabular data either as the predictor itself, on serialized rows
\citep{hegselmann2023tabllm,dinh2022lift}, or to engineer features and rules for a model that is
then trained or retrained: CAAFE generates features from dataset descriptions
\citep{hollmann2023caafe}, FeatLLM turns LLM-written rules into features for few-shot learning
\citep{han2024featllm}, and OCTree refines feature-generation rules with decision-tree feedback
\citep{nam2024octree}. Closest to us, DeLTa has an LLM rewrite a tree ensemble's rules into one that
corrects the ensemble's outputs \citep{ye2025delta}. These methods target accuracy on public
benchmarks, where the predictor can be rebuilt freely. Our setting differs in three ways: the model is
deployed and frozen, so the rules act as a reversible overlay at its existing cut-off; generation is
anchored on the features the model underweights rather than on refining its existing tree rules; and the
training data are selection-biased, so the gain must hold on a population the model was not trained
on.

As the rule proposer, a single LLM produces high-variance outcomes not predictable from in-distribution
signal: the backbone that performs best on the real-world generalization test is among the worst on
the BAF variants, and vice versa (Section~\ref{sec:results}). We therefore prompt three
architecturally diverse backbones \citep{wei2022cot,deepseek2025r1,qwen2026qwen35,gemma2026gemma4} across
independent shots and average all retained shots in logit space \citep{dietterich2000ensemble}, with
no dataset-specific selection.

\subsection*{Contributions}

\begin{enumerate}[nosep,leftmargin=1.2em]
  \item \textbf{Patching a frozen, deployed model.} LLM rules steered through the detector's
    SHAP-underweighted dimensions, calibrated into a reversible overlay that runs at the production
    cut-off with every parameter fixed before deployment.
  \item \textbf{Evidence under selection bias and designed shift.} On an unbiased cohort, continued
    training regresses while the overlay gains recall at equal precision; on seven bias-shift
    variants it raises precision at equal recall on every variant. Pooling three LLMs bounds the high,
    dataset-dependent variance of single-LLM benefit without oracle selection
    (Section~\ref{sec:results}).
\end{enumerate}

\section{Method}
\label{sec:method}

\sysname\ adds a calibrated, LLM-generated rule overlay to a frozen gradient-boosted classifier.
Figure~\ref{fig:framework} summarizes the pipeline, from a model-derived prompt context through
LLM proposal, empirical screening, band-conditioned calibration, and a reversible overlay (top),
and how the overlay is deployed beside the production score (bottom).

\definecolor{colCtx}{RGB}{214,229,252}    %
\definecolor{colLLM}{RGB}{252,238,210}    %
\definecolor{colScreen}{RGB}{210,245,218} %
\definecolor{colCalib}{RGB}{237,226,250}  %
\definecolor{colOv}{RGB}{215,242,230}     %

\begin{figure*}[t]
\centering
\resizebox{\textwidth}{!}{%
\begin{tikzpicture}[
  >=Latex,
  node distance=0.45cm,
  box/.style={draw, rounded corners=3pt, align=center, inner sep=5pt,
              font=\footnotesize},
  mbox/.style={draw, rounded corners=2pt, align=center, inner sep=2pt,
               font=\scriptsize, fill=colLLM},
  arr/.style={->, thick},
  darr/.style={->, dashed, semithick, gray!65!black},
  anarr/.style={->, thin, gray!55!black, shorten >=2pt, shorten <=2pt},
]

%
\node[box, fill=colCtx, text width=1.95cm, minimum height=2.3cm] (ctx) {%
  \textbf{Prompt}\\
  \textbf{Context}\\[3pt]
  \scriptsize domain primer\\[2pt]
  model digest\\
  (splits, gain,\\
  paths)\\[2pt]
  SHAP gap $+$\\
  lift tables%
};
\node[font=\scriptsize, align=center, text width=2.05cm, text=gray!68!black,
      below=0.35cm of ctx] (ctxa)
  {\textit{summaries of the}\\
   \textit{frozen GBM;}\\
   \textit{model unchanged}};
\draw[anarr] (ctx.south) -- (ctxa.north);

%
\node[mbox, right=0.8cm of ctx, yshift=0.95cm, text width=1.9cm] (llm1) {%
  \textbf{LLM-CoT}\\
  chain-of-thought%
};
\node[mbox, below=0.22cm of llm1, text width=1.9cm] (llm2) {%
  \textbf{LLM-SlW}\\
  sliding-window%
};
\node[mbox, below=0.22cm of llm2, text width=1.9cm] (llm3) {%
  \textbf{LLM-MoE}\\
  sparse MoE%
};
\begin{pgfonlayer}{background}
  \node[draw, dashed, rounded corners=5pt, fill=colLLM!30,
        fit=(llm1)(llm2)(llm3), inner sep=7pt] (llmg) {};
\end{pgfonlayer}
\node[font={\scriptsize\bfseries}, anchor=south west]
  at ($(llmg.north west) + (3pt,2pt)$)
  {3 open-weight LLMs $\times$ 10 independent shots};

%
\node[box, fill=colScreen, right=0.85cm of llmg,
      text width=2.05cm, minimum height=2.3cm] (scr) {%
  \textbf{Support}\\
  \textbf{Screen}\\[3pt]
  \scriptsize compiles, fires,\\
  lift $>1$;\\[2pt]
  per band: $\geq 3$\\
  fires, lift $>1$%
};

%
\node[box, fill=colCalib, right=0.75cm of scr,
      text width=2.05cm, minimum height=2.3cm] (cal) {%
  \textbf{Band-Cond.}\\
  \textbf{Log-Lift}\\
  \textbf{Calibration}\\[3pt]
  \scriptsize log-lift residual\\
  per alert-rate\\
  band%
};

%
\node[box, fill=colOv, right=0.75cm of cal,
      text width=2.05cm, minimum height=2.3cm] (ov) {%
  \textbf{Logit-Space}\\
  \textbf{Averaging}\\
  $+$\,\textbf{Overlay}\\[3pt]
  \scriptsize uniform avg.\\
  across retained\\
  shots \& models;\\
  \emph{additive}\\
  overlay%
};
\node[font=\scriptsize, align=center, text width=2.2cm, text=gray!68!black,
      below=0.35cm of ov] (ova)
  {\textit{reversible: base GBM}\\
   \textit{never modified;}\\
   \textit{disable w/o redeploy}};
\draw[anarr] (ov.south) -- (ova.north);

%
\node[box, fill=white, right=0.75cm of ov,
      text width=1.55cm, minimum height=2.2cm] (out) {%
  \textbf{Final}\\
  \textbf{Score}\\[3pt]
  \scriptsize frozen GBM\\
  $+$ overlay%
};

%
\draw[arr] (ctx.east) -- (llmg.west);
\draw[arr] (llmg.east) --
  node[above, font=\scriptsize, fill=white, inner sep=1.5pt]{candidate rules}
  (scr.west);
\draw[arr] (scr.east) -- (cal.west);
\draw[arr] (cal.east) -- (ov.west);
\draw[arr] (ov.east) -- (out.west);

\end{tikzpicture}%
}
\par\medskip
\resizebox{\textwidth}{!}{%
\begin{tikzpicture}[
  >=Latex,
  box/.style={draw, rounded corners=3pt, align=center, inner sep=3.5pt, font=\footnotesize, minimum height=1.05cm},
  arr/.style={->, thick},
  oarr/.style={->, thick, teal!50!black},
]
\node[font=\scriptsize\bfseries, anchor=west, text=gray!60!black] at (-0.9, 1.55) {PRODUCTION (unchanged)};
\node[font=\scriptsize\bfseries, anchor=west, text=teal!50!black] at (-0.9, -1.75) {OVERLAY (same cut-off $t_b$)};
\node[box, fill=white, text width=1.65cm] (x) at (0,0) {\textbf{Transaction}\\$\mathbf{x}$};
\node[box, fill=colCtx, text width=2.6cm] (f) at (3.1,0.85) {\textbf{Frozen model} $f$\\\scriptsize production score $s(\mathbf{x})$};
\node[box, fill=white, text width=2.3cm] (pa) at (18.2,0.85) {\textbf{Production alert}\\\scriptsize $s(\mathbf{x})\ge t_b$};
\node[box, fill=colCalib, text width=2.6cm] (r) at (3.1,-0.85) {\textbf{Rule pools} $\mathcal{K}_{j,b}$\\\scriptsize $R_k(\mathbf{x})$, $\log\mathrm{lift}_k$};
\node[box, fill=colOv, text width=2.85cm] (u) at (6.9,-0.85) {\textbf{Base percentile}\\\scriptsize $p(\mathbf{x})=F_A\bigl(s(\mathbf{x})\bigr)$};
\node[box, fill=colOv, text width=3.9cm] (z) at (10.85,-0.85) {\textbf{Overlay score}\\\scriptsize $z_b(\mathbf{x})=\tfrac{1}{|S|}\sum_{j}\operatorname{logit}p^*_{j,b}(\mathbf{x})$};
\node[box, fill=colOv, text width=3.35cm] (st) at (14.85,-0.85) {\textbf{Scaled overlay score}\\\scriptsize $\tilde s_b=\max\{t: z_b\ge\tau_b(t)\}$};
\node[box, fill=colOv, text width=2.3cm] (oa) at (18.2,-0.85) {\textbf{Overlay alert}\\\scriptsize $\tilde s_b(\mathbf{x})\ge t_b$};
\draw[arr] (x.east) -- ++(0.25,0) |- (f.west);
\draw[arr] (x.east) -- ++(0.25,0) |- (r.west);
\draw[arr] (f.east) -- (pa.west);
\draw[oarr] ($(f.east)+(0.45,0)$) |- (u.west);
\draw[oarr] (u.east) -- (z.west);
\draw[oarr] (r.south) -- ++(0,-0.3) -| (z.south);
\draw[oarr] (z.east) -- (st.west);
\draw[oarr] (st.east) -- (oa.west);
\node[font=\scriptsize, text=gray!65!black, anchor=west] at (5.4,-2.05)
  {\textit{Fitted once on the adaptation window $A$, then frozen:} percentile table $F_A$; lifts; cut-off table $\tau_b(t)$, which alerts as many rows of $A$ as $s\ge t$.};
\node[font=\scriptsize, text=gray!65!black, anchor=west] at (5.4,-2.42)
  {\textit{Rules off} $\Rightarrow \tilde s_b = s$\textit{: the overlay is switched off without redeploying $f$.}};
\end{tikzpicture}}
\caption{\textbf{Top: framework overview (left to right).}
  A prompt context built from frozen base-model outputs drives three open-weight LLMs; emitted rules
  are support-screened and calibrated as band-conditioned log-lift
  residuals, giving an \emph{additive, reversible} overlay on the frozen base score.
  \textbf{Bottom: deployment.} The production score and cut-off are unchanged; the overlay is mapped
  back to the same scale by a cut-off table fitted on the adaptation window and alerts at the same
  cut-off $t_b$ (Section~\ref{sec:method:ensemble}).}
\label{fig:framework}
\end{figure*}

\subsection{Problem Setup and Frozen-Model Constraint}
\label{sec:method:setup}

Let $\mathbf{x}\in\mathcal{X}$ be a transaction feature vector and $y\in\{0,1\}$ the fraud label.
A tree-ensemble base scorer $f$ (constructed in Section~\ref{sec:setup}) is trained once on
historical data and frozen thereafter. Its production score $s(\mathbf{x})\in[0,1000]$ is alerted
at a cut-off $t$, i.e.\ when $s(\mathbf{x})\ge t$. Our framework never modifies $f$, $s$, or $t$.
Because $f$ is never retrained, the generalization-gap failure mode described above is avoided by
construction (Section~\ref{sec:results}).
We call the labeled slice used to generate, screen, and calibrate the overlay the
\emph{adaptation window} $A$; it is not a held-out evaluation set. Every fitted quantity---rules,
lifts, and the lookup tables below---comes from $A$ and is frozen: no out-of-time (OOT) labels or
score distributions are used at any stage. The overlay's base is $p(\mathbf{x})=F_A(s(\mathbf{x}))$,
the score's percentile under $A$'s empirical distribution (a fixed lookup table).

\subsection{Model-Aware LLM Rule Proposal}
\label{sec:method:adv}

For each dataset the framework runs \textbf{3 LLMs $\times$ 10 independent generation shots}: each
shot is a single generation pass with the same prompt (Appendix~\ref{sec:app_prompt}), yielding a
noisy candidate dictionary of diverse predicates.

The prompt has three parts: (i) a domain primer describing the fraud-detection context and feature
encodings; (ii) a digest of the frozen model's trees---per-feature split counts and gain, frequent
feature sets on high-gain paths, and low-value leaves; and (iii) a SHAP attribution gap on the
adaptation window: features are split at the median of their mean SHAP over fraud transactions into
those the model relies on and those it underweights, with decile (or category) fraud-lift tables for
the underweighted ones. The LLM sees these summaries, never the model itself. It reasons about how
fraud would present in the \emph{underweighted} dimensions---behaviours a detector anchored on the
first group would not flag---and every candidate's \emph{primary anchor} must come from that set;
strongly-modeled features may appear only as secondary filters narrowing a pattern to a behavioural
archetype. Each candidate must cite supporting lift from the tables. Candidates are conjunctions of
threshold conditions $(\text{feature}\;\circ\;\theta)$, $\circ\in\{<,\le,>,\ge,=\}$, or missing /
non-missing tests, joined by $\wedge$ (no disjunctions), with thresholds given as training-set
percentiles.

\paragraph{Why an LLM rather than a rule miner?}
The adaptation window's per-band strata are small, so a tree learner or itemset miner fits only
patterns \emph{already present}, while the LLM proposes archetypes from domain priors; its output is
also a named hypothesis with a rationale (Section~\ref{sec:abl:patterns}). We do not compare
empirically against rule-mining baselines and flag this as the primary gap.

The three backbones (exact identifiers in Section~\ref{sec:setup}) span chain-of-thought reasoning,
sliding-window attention, and sparse mixture-of-experts; no single one wins uniformly across
datasets (Section~\ref{sec:results}), motivating the all-source ensemble.

\subsection{Empirical Support Screening}
\label{sec:method:screen}

Predicates are compiled with their percentile thresholds taken from the training data; any that fail
to parse or compile are dropped. On the adaptation window, a predicate is kept if it fires, its fraud
rate exceeds the window's base rate (lift $>1$), and it is not a duplicate of an already kept
predicate (canonical form); on BAF, predicates built only from strongly-modeled features are also
dropped. It then enters the pool $\mathcal{K}_b$ of every band $b$ in whose score
zone $Z_b$ (Section~\ref{sec:method:nb}) it fires on at least three transactions with lift $>1$.
Shots left with no rule in any band are discarded (the \textbf{compilable-shot filter}); retained
shot counts are reported in Section~\ref{sec:setup}. Screening checks empirical support only; it does
not prove semantic correctness or guarantee out-of-time transfer.

\subsection{Band-Conditioned Log-Lift Residual Overlay}
\label{sec:method:nb}

We operate at a set of production cut-offs $t_b$ (alert-rate bands $b$). For each band the
framework fits an independent rule pool $\mathcal{K}_b$ and lift set on the adaptation window
\emph{restricted to band $b$'s score zone} $Z_b=\{\mathbf{x}\in A: s(\mathbf{x})\ge t_b\}$.
Let $R_k(\mathbf{x})\in\{0,1\}$ indicate whether rule $k$ fires on $\mathbf{x}$.
For a rule $k\in\mathcal{K}_b$, its calibrated fraud lift is the band-conditioned precision ratio
\begin{equation}
  \mathrm{lift}_k \;=\;
    \frac{P\!\left(y=1 \mid R_k=1,\; \mathbf{x}\in Z_b\right)}
         {P\!\left(y=1 \mid \mathbf{x}\in Z_b\right)},
  \label{eq:lift}
\end{equation}
estimated by a decile-standardized average: $Z_b$ is split into base-score deciles, and the
within-decile precision ratios are averaged weighted by rule-firing support. This controls for
score-dependent base rates and targets---rather than exactly equals---Eq.~\eqref{eq:lift}.
The overlay for cut-off $t_b$ applies $\mathcal{K}_b$ to every transaction:
\begin{equation}
  \boxed{%
  \begin{aligned}
    p^*_b(\mathbf{x}) ={}& \sigma\Bigl(\operatorname{logit}\bigl(p(\mathbf{x})\bigr) \\
    &{}+ \!\!\sum_{k\in\mathcal{K}_b}\!\! \log(\mathrm{lift}_k)\,R_k(\mathbf{x})\Bigr)
  \end{aligned}}
  \label{eq:overlay}
\end{equation}
where $\sigma(\cdot)$ is the sigmoid and $p(\mathbf{x})=F_A(s(\mathbf{x}))$
(Section~\ref{sec:method:setup}). A firing rule with $\mathrm{lift}_k>1$ raises the score, one
below~1 lowers it ($\log$ is natural, matching the logit scale).

\textbf{Per-band calibration} is load-bearing: estimated globally, a rule's coefficient is diluted
by transactions outside its active zone, and D-Industrial avg$\Delta$R@P falls from $+2.26$ to
$-2.41\%$ (Appendix~\ref{sec:app_pipeabl}).

\subsection{Multi-Shot Averaging, Scaled Score, and Evaluation}
\label{sec:method:ensemble}

Each retained shot $j$ (across all three LLMs) produces $p^*_{j,b}(\mathbf{x})$ via
Eq.~\eqref{eq:overlay}. Stochastic generation produces heterogeneous candidate pools across shots;
the ensemble combines all retained shots $S$ by \textbf{uniform averaging in logit space},
$z_b(\mathbf{x}) = |S|^{-1}\sum_{j\in S}\operatorname{logit}\bigl(p^*_{j,b}(\mathbf{x})\bigr)$.
Every shot receives equal weight and the same combination is applied to all datasets;
logit-space averaging matches the additive parameterization of Eq.~\eqref{eq:overlay} and is a
standard variance-reduction strategy \citep{dietterich2000ensemble}.

\paragraph{Scaled overlay score and deployment.}
To run at the existing cut-off, $z_b$ is mapped back to the production scale by a cut-off table
fitted on $A$: $\tau_b(t)$ is the value at which $z_b$ alerts as many transactions of $A$ as
$s\ge t$, for every $t$ on a 0.1-point grid, and
\begin{equation}
  \tilde s_b(\mathbf{x}) = \max\{\,t : z_b(\mathbf{x}) \ge \tau_b(t)\,\}.
  \label{eq:scaled}
\end{equation}
The overlay alerts when $\tilde s_b(\mathbf{x})\ge t_b$, the production cut-off
(Figure~\ref{fig:framework}, bottom). Serving needs only $F_A$, the rules with their lifts, and
$\tau_b$; nothing is refitted on new traffic, and with the rules switched off $\tilde s_b=s$.

\paragraph{Evaluation metric.}
On OOT data we compare production ($s\ge t$) and the overlay ($\tilde s_b\ge t$) through their
precision--recall points over the same threshold grid. Let $P_b$, $R_b$ be production's precision
and recall at its cut-off $t_b$. \textbf{$\Delta$R@P} is the overlay's highest recall at any threshold
with precision $\ge P_b$, minus production's own highest recall at precision $\ge P_b$;
\textbf{$\Delta$P@R} is defined symmetrically at recall $\ge R_b$. Subtracting production's own best
removes what it would gain by merely moving its threshold; no interpolation is used, and a level the
overlay never reaches counts as recall (precision) zero. We report averages over a split's
cut-offs, \textbf{avg$\Delta$R@P} and \textbf{avg$\Delta$P@R}, as absolute differences in \%; the
precision and recall levels they average over are given in Section~\ref{sec:setup}, and every
operating point in Appendix~\ref{sec:app_perband}. This pair is the sole comparison metric; we do
not use extra-TP counts or PR-AUC/AUC as comparison axes.

\section{Experimental Setup}\label{sec:setup}

\paragraph{Datasets: generalization and robustness to variants.}
We evaluate on two datasets that test the framework's two target capabilities
(Table~\ref{tab:datasets}): generalizing to an unseen population, and staying robust across designed
distribution-shift variants.
\textbf{D-Industrial} detects unauthorized-transaction fraud at a real-world e-commerce platform,
trained on sampled data: training covers approved traffic only, richer in fraud signal but
selection-biased, while evaluation uses a small unbiased cohort approved regardless of score
(Section~\ref{sec:intro}). This is where continued training generalizes worst.
\textbf{BAF}~\citep{jesus2022baf} is a public bank-account benchmark with seven OOT splits of
designed bias-shift variants and no train/evaluation population mismatch: gradient adaptation is
well-posed there, so it tests whether the overlay stays robust when retraining is a viable
alternative.

\begin{table*}[t]
  \centering
  \footnotesize
  \setlength{\tabcolsep}{5pt}
  \caption{Datasets and base-model characteristics. Both use out-of-time evaluation; D-Industrial is
    sampled from real-world transaction data and is not public, BAF is a public benchmark.
    Blend $\alpha$ is the LightGBM weight in the frozen base's rank blend with a random forest,
    swept on both datasets and selected by validation PR-AUC (sample-weighted on D-Industrial); $\alpha{=}1$ is the sweep's outcome,
    not an untuned model.}
  \label{tab:datasets}
  \begin{tabular}{@{}l p{2.7cm} r r c p{5.3cm}@{}}
    \toprule
    Dataset & Shift type & \#feat & \#trees (LGB\,/\,RF) & Blend $\alpha$ & Split design \\
    \midrule
    D-Industrial
      & selection bias; unbiased eval $\neq$ biased train
      & 500 & 2{,}219\,/\,500 & $0.75$
      & Temporal, six months: train Oct--Nov 2025 (biased, approved-only); validate Dec 2025 and
        evaluate Jan 1--Mar 7, 2026, both from the unbiased cohort \\[2pt]
    BAF
      & public; designed bias-shift variants
      & 30 & 892\,/\,200 & $1.00$
      & Temporal: train months 0--4; evaluate held-out months 5--7 (OOT-1, OOT-2) plus five
        designed bias-shift variants (Var-I--V) \\
    \bottomrule
  \end{tabular}
\end{table*}

\paragraph{Baselines.}
\textbf{Frozen base} is the scorer under test, never modified; all $\Delta$ metrics are relative to
it. It is itself ensemble-tuned: a rank blend
$\alpha\,\mathrm{rank}(\mathrm{LGB})+(1{-}\alpha)\,\mathrm{rank}(\mathrm{RF})$
\citep{ke2017lightgbm}, ranks taken against the adaptation window, with $\alpha$ swept over a grid
and kept at the value maximizing validation PR-AUC (Table~\ref{tab:datasets}), so overlay gains are
measured over an already ensemble-tuned baseline. \textbf{CT-100} adds 100 boosting rounds on the
adaptation window (a straightforward continued-training baseline). \textbf{CT-Tuned} applies one
boosting round at a lower learning rate (lr$=$0.03), a single round-count/learning-rate sensitivity
point evaluated on D-Industrial only. \textbf{Retrain} fully retrains on the combined
historical-plus-adaptation data. Each is scored on the same scale and evaluated like the overlay,
with its cut-off table fitted on the adaptation window.

\paragraph{Operating cut-offs.}
All methods are evaluated at the production cut-offs $t_b$ on the 0--1000 score scale:
D-Industrial $\{998,\dots,993\}$ (6) and BAF $\{999,995,990,985,980,975,970\}$ (7). At these
cut-offs production's OOT precision and recall span 30.3--56.5\% and 23.5--62.9\% on D-Industrial,
and 10.1--56.2\% and 3.1--50.1\% across the BAF splits; avg$\Delta$R@P and avg$\Delta$P@R average
over these precision and recall levels respectively.

\paragraph{LLM backbones and active shot counts.}
Three architecturally diverse open-weight backbones each generate 10 independent shots per dataset:
\textbf{LLM-CoT} = \texttt{DeepSeek-R1-\allowbreak Distill-Qwen-\allowbreak 32B-4bit}~\citep{deepseek2025r1}
(chain-of-thought); \textbf{LLM-SlW} = \texttt{Gemma4-31B-it-8bit}~\citep{gemma2026gemma4}
(sliding-window attention); \textbf{LLM-MoE} = \texttt{Qwen3.5-35B-A3B-4bit}~\citep{qwen2026qwen35}
(mixture-of-experts). Shots left with no rule after screening are discarded; active counts after
filtering (LLM-CoT\,/\,LLM-SlW\,/\,LLM-MoE) are D-Industrial \mbox{9\,/\,8\,/\,8} and BAF
\mbox{10\,/\,6\,/\,9} (19/14/17 total). The all-3 ensemble averages uniformly over all retained
shots. Generation ran locally on one Apple M4 Max workstation (128\,GB unified memory); no hosted
API was used, so transaction data never leaves the environment.

%
\section{Results}\label{sec:results}

\subsection{Continued Training vs.\ the Overlay}\label{sec:results:conttrain}

We compare the all-3 ensemble overlay with three continued-training variants---CT-Tuned, CT-100, and
full Retrain---under the same avg$\Delta$R@P / avg$\Delta$P@R metric (full figures in
Appendix~\ref{sec:app_conttrain}, Table~\ref{tab:conttrain}).

\begin{table*}[t]
  \centering
  \footnotesize
  \renewcommand{\arraystretch}{0.95}
  \caption{Per-backbone and ensemble results: \textbf{avg$\Delta$R@P} / avg$\Delta$P@R (absolute
    differences, \%) relative to the frozen base at its production cut-offs (Section~\ref{sec:method:ensemble}).
    \emph{Anchors}: production's OOT precision and recall at the split's cut-offs, the fixed levels at
    which $\Delta$R@P and $\Delta$P@R are taken. \textbf{Ensemble} is the primary result (bold where
    positive). Var-IV shares Var-III's population. CoT, SlW, MoE: Section~\ref{sec:setup}.}
  \label{tab:main}
  \resizebox{\textwidth}{!}{%
  \begin{tabular}{@{}ll cc rrrr@{}}
    \toprule
    & & \multicolumn{2}{c}{Anchors (production)} & & & & \\
    \cmidrule(lr){3-4}
    Dataset & Split & Precision & Recall & CoT & SlW & MoE & \textbf{Ensemble (All 3)} \\
    \midrule
    D-Industrial & OOT & 30.3--56.5\% & 23.5--62.9\% & $-0.28\%$ / $-0.15\%$ & $+2.27\%$ / $+1.56\%$ & $+3.88\%$ / $+3.09\%$ & $\mathbf{+2.26\%}$ / $\mathbf{+1.60\%}$ \\
    \midrule
    BAF & OOT-1 & 17.3--53.2\% & 4.1--41.4\% & $-0.12\%$ / $+0.12\%$ & $-0.21\%$ / $+0.06\%$ & $-0.18\%$ / $-0.09\%$ & $-0.10\%$ / $\mathbf{+0.07\%}$ \\
     & OOT-2 & 19.1--52.8\% & 3.9--41.2\% & $-0.08\%$ / $+0.09\%$ & $-0.41\%$ / $-0.17\%$ & $-0.26\%$ / $-0.14\%$ & $0.00\%$ / $\mathbf{+0.06\%}$ \\
     & Var-I & 18.0--56.2\% & 3.1--35.3\% & $+0.21\%$ / $+0.12\%$ & $-0.07\%$ / $+0.17\%$ & $-0.22\%$ / $+0.04\%$ & $\mathbf{+0.04\%}$ / $\mathbf{+0.08\%}$ \\
     & Var-II & 14.6--51.8\% & 6.4--50.1\% & $+0.37\%$ / $+0.25\%$ & $+0.20\%$ / $+0.13\%$ & $+0.19\%$ / $+0.13\%$ & $\mathbf{+0.27\%}$ / $\mathbf{+0.30\%}$ \\
     & Var-III & 10.1--36.6\% & 4.5--43.1\% & $+0.23\%$ / $+0.26\%$ & $-0.22\%$ / $+0.18\%$ & $-0.07\%$ / $+0.41\%$ & $\mathbf{+0.04\%}$ / $\mathbf{+0.28\%}$ \\
     & Var-IV & 10.1--36.6\% & 4.5--43.1\% & $+0.23\%$ / $+0.26\%$ & $-0.22\%$ / $+0.18\%$ & $-0.07\%$ / $+0.41\%$ & $\mathbf{+0.04\%}$ / $\mathbf{+0.28\%}$ \\
     & Var-V & 10.8--40.4\% & 5.1--45.6\% & $+0.10\%$ / $+0.01\%$ & $-0.51\%$ / $-0.13\%$ & $-0.36\%$ / $-0.28\%$ & $-0.13\%$ / $\mathbf{+0.02\%}$ \\
    \bottomrule
  \end{tabular}%
  }
\end{table*}

\paragraph{D-Industrial: the generalization test.}
The \emph{fair} comparator here is a carefully tuned continued-training update, not a naive one. Even
CT-Tuned (one boosting round, lr$=$0.03) regresses, at avg$\Delta$R@P $-7.69\%$ and avg$\Delta$P@R
$-2.85\%$; the naive CT-100 and full Retrain never reach production's precision at any cut-off (avg$\Delta$R@P $-46.39\%$: no recall at those precision levels) and lose $34.80\%$ and $41.73\%$ of precision at equal recall. The overlay instead
\emph{gains} avg$\Delta$R@P $+2.26\%$ (95\% CI $[+1.84\%,+2.90\%]$, paired block bootstrap
over 66 calendar-day blocks, 1000 replicates) and avg$\Delta$P@R $+1.60\%$
($[+1.28\%,+2.04\%]$), positive at every cut-off over the OOT period and on average in each month
(Appendix~\ref{sec:app_month}), with no redeployment. The cause
is structural, not under-tuning: every training example carries the same selection bias, so no
learning rate or round count turns the biased gradient direction into an unbiased
one~\citep{quinonero2009shift}; a frozen base sidesteps the mechanism entirely.

\paragraph{BAF: a robustness floor.}
BAF tests robustness to designed variants rather than generalization across a population gap, so
gradient updates are well-posed and the overlay should not be expected to lead---and it does not.
What it shows is that the overlay does no harm: avg$\Delta$P@R is positive on all seven splits
($[+0.02\%,+0.30\%]$) and avg$\Delta$R@P stays near neutral ($[-0.13\%,+0.27\%]$;
negative on OOT-1 and Var-V). CT-100 is no better ($[-0.39\%,-0.04\%]$ / $[-0.29\%,+0.32\%]$); Retrain raises precision at equal recall more ($[+0.08\%,+0.71\%]$) but loses recall at equal precision on six splits (down to $-0.75\%$), and both need a redeploy. The claim is a
robustness floor, not an improvement.

\subsection{Proposal-Source Variance and Ensemble Stabilization}\label{sec:results:main}


Table~\ref{tab:main} reports both metrics per backbone and for the all-3 ensemble. The central
finding is \emph{high single-LLM variance}: the regimes invert the ranking---LLM-MoE is best
industrially ($+3.88\%$ avg$\Delta$R@P) yet negative on six of seven BAF splits, while
LLM-CoT is slightly negative on D-Industrial ($-0.28\%$) yet the best single source on
\emph{all seven}. No in-distribution signal separates them.

Ensembling all three resolves this without an oracle: avg$\Delta$P@R is positive on \emph{all eight}
splits and avg$\Delta$R@P on five (the exceptions within $0.13\%$ of neutral), trading part
of LLM-MoE's industrial peak ($+2.26\%$ vs $+3.88\%$) to avoid its BAF degradation.
The gain also needs the rest of the pipeline (Appendix~\ref{sec:app_pipeabl}): without the log-lift
overlay or the per-shot ensemble, the same LLM rules fall below the frozen base on every split
(avg$\Delta$R@P $-2.5$ to $-46\%$ on D-Industrial, $-17$ to $-33\%$ on BAF), and global instead of
per-band calibration lowers avg$\Delta$R@P on every split.

\section{Analysis of LLM-Proposed Rules}
\label{sec:ablation}

\subsection{LLM Source and Proposal Behaviour}

Appendix~\ref{sec:app_srcabl} (Table~\ref{tab:srcabl_full}) covers all seven subsets. The
inversion carries over to pairs: LLM-SlW$+$LLM-MoE, the best pair on D-Industrial
($+2.80\%$ avg$\Delta$R@P), is worst or near-worst on every BAF split; LLM-CoT pairs invert this.

The backbones also differ in how they follow the output contract (Appendix~\ref{sec:app_prompt}). On
D-Industrial, 44\% of LLM-CoT's returned predicates name features that do not exist in the
model---half of them near-misses of real names---against 2\% or less for the other two, and
62\% of LLM-SlW's are repeats within a shot; on both datasets, 69--73\% of LLM-MoE's pass the support
screen. The retained rule sets barely overlap: only 14\% (D-Industrial) and 16\% (BAF) of distinct
retained predicates come from more than one shot, and LLM-CoT shares just 2 of its 22 D-Industrial
rules with each other backbone. Each backbone thus contributes largely different patterns---the
variance the ensemble absorbs.

\subsection{Representative Risk Patterns}
\label{sec:abl:patterns}

Each retained predicate carries the LLM's name and rationale beside its measured lift. The LLMs
largely follow the steering instruction: the first condition of 97\% (D-Industrial) and 99\% (BAF) of
retained predicates is an underweighted feature, with strongly-modeled features as secondary filters in
9\% and 33\%. Rules typically combine individually weak evidence into a selective conjunction. On
D-Industrial, LLM-CoT cites a $1.09\times$ lift for failed payments at one shipping address and proposes
\texttt{\scriptsize failed\_txn(addr,14d)>p90 $\wedge$ buyer\_seg>p80}, reasoning that repeated failed
payments to one address test card authorization; on the adaptation window the rule reaches
$71.6\times$ lift on 6{,}249 transactions. On BAF, LLM-MoE combines \texttt{\scriptsize session\_len>p90}
($1.44\times$) and \texttt{\scriptsize free\_email} ($1.22\times$) with \texttt{\scriptsize housing=0}
into a $6.7\times$ rule (862 transactions). Calibration lift does not predict which generalize: the
highest-lift D-Industrial predicates come from LLM-CoT, the weakest single source there.

\section{Conclusion}\label{sec:conclusion}

We patched a frozen fraud detector's blind spots by steering LLM proposals through its
underweighted dimensions into a reversible log-lift overlay that runs at the production cut-off. It
generalizes where continued training regresses, stays robust
where retraining offers no edge, and reports each amendment as an auditable predicate; pooling
backbones bounds the worst case.

\section*{Limitations}

\paragraph{Scope of the recall advantage.}
The recall advantage is demonstrated under a generalization gap between training and evaluation
populations. On BAF, where no such gap exists, neither method dominates on recall and we claim only
that the overlay does no harm.

\paragraph{Continued-training baselines.}
We evaluate two continued-training configurations on D-Industrial (CT-100 and CT-Tuned); both trail
the frozen baseline, but two configurations do not span the hyperparameter space. Importance-weighted
or propensity-adjusted continued training is the principled remedy; we do not evaluate it, as the
selection mechanism is a separate model whose scores are unavailable to the adaptation procedure.

\paragraph{Inference cost.}
The three-model ensemble requires roughly $3\times$ the LLM inference budget of a single-source
overlay; a single-source fallback is feasible but sacrifices variance reduction.

\paragraph{Offline evaluation.}
All results are offline evaluations against historical labels, with every parameter fixed on the
adaptation window; no live A/B test has been conducted.

\section*{Ethics Statement}

\paragraph{Data and privacy.}
Experiments on D-Industrial use historical transaction records accessed under the platform's
internal data-governance program. No human-subject data were collected, no directly identifying
attributes are used as model features, and feature names are paraphrased throughout. BAF is a public
benchmark released for fraud-detection research.

\paragraph{Risk of harm.}
Fraud decisions carry asymmetric costs: a false alert burdens a legitimate customer, a missed fraud
harms the platform and its users. The overlay alerts at the production cut-off and never blocks
autonomously; deployments should retain human review and an appeals path.

%
\bibliography{custom}

%
\appendix

\section{Continued-Training Comparison}
\label{sec:app_conttrain}

Table~\ref{tab:conttrain} gives the full continued-training comparison summarized in
Section~\ref{sec:results:conttrain}: the all-3 ensemble overlay against three continued-training
variants (CT-100, CT-Tuned, full Retrain) under the same avg$\Delta$R@P / avg$\Delta$P@R metric,
across both datasets. Each method's score is mapped to the production scale by its own cut-off table
fitted on the adaptation window (Section~\ref{sec:method:ensemble}).

\begin{table*}[t]
  \centering
  \footnotesize
  \setlength{\tabcolsep}{3pt}
  \caption{Comparison with continued training: \textbf{avg$\Delta$R@P} / avg$\Delta$P@R (\%)
    relative to the frozen base at its production cut-offs. BAF figures are ranges over the 7 splits.
    On D-Industrial, CT-100 and Retrain never reach production's precision at any cut-off, so
    their recall at those precision levels is zero. CT-Tuned (1 round, lr$=$0.03) evaluated on
    D-Industrial only. \textbf{Bold column} = our framework (all-3 ensemble).}
  \label{tab:conttrain}
  \resizebox{\textwidth}{!}{%
  \begin{tabular}{@{}l rr rr rr rr@{}}
    \toprule
    & \multicolumn{2}{c}{CT-100} & \multicolumn{2}{c}{CT-Tuned} & \multicolumn{2}{c}{Retrain}
    & \multicolumn{2}{c}{\textbf{Ours}} \\
    \cmidrule(lr){2-3}\cmidrule(lr){4-5}\cmidrule(lr){6-7}\cmidrule(lr){8-9}
    Dataset (split) & \textbf{$\Delta$R@P} & $\Delta$P@R
                    & \textbf{$\Delta$R@P} & $\Delta$P@R
                    & \textbf{$\Delta$R@P} & $\Delta$P@R
                    & \textbf{$\Delta$R@P} & $\Delta$P@R \\
    \midrule
    D-Industrial OOT
      & $-46.39\%$ & $-34.80\%$
      & $-7.69\%$ & $-2.85\%$
      & $-46.39\%$ & $-41.73\%$
      & $\mathbf{+2.26\%}$ & $\mathbf{+1.60\%}$ \\
    BAF (7 splits, range)
      & $[-0.39\%,-0.04\%]$ & $[-0.29\%,+0.32\%]$
      & \multicolumn{2}{c}{---}
      & $[-0.75\%,+0.09\%]$ & $[+0.08\%,+0.71\%]$
      & $[-0.13\%,+0.27\%]$ & $[+0.02\%,+0.30\%]$ \\
    \bottomrule
  \end{tabular}%
  }
\end{table*}

\section{Prompt Design}
\label{sec:app_prompt}

Each generation shot sends a single prompt to the LLM; there is no second LLM pass.
Listing~\ref{lst:prompt} reproduces the template behind all reported results (D-Industrial and BAF,
all three backbones), as recovered from the logged LLM calls. Four placeholders are filled per
dataset:
\begin{itemize}[nosep,leftmargin=1.1em,topsep=2pt]
  \item \texttt{\{domain primer\}}: a description of the fraud-detection task and of the dataset's
    features and their encodings. It is dataset-specific and not reproduced; for D-Industrial it
    describes proprietary data.
  \item \texttt{\{model digest\}}: a summary of the frozen model's trees---split counts, total gain and
    split depths for the features with the most gain, the names of the remaining features, the
    feature sets most frequent on high-gain leaf paths, and a few low-value leaves.
  \item \texttt{\{SHAP gap\}}: features split at the median of their mean SHAP over fraud transactions
    on the adaptation window, into those above it (``forbidden'', listed with their mean SHAP) and
    those at or below it (``target''), each target feature with a decile or category table of counts,
    fraud rate, and lift over the window's base rate. The template's ``SHAP \texttt{<=} 0'' refers to
    this target set.
  \item \texttt{\{number of rules\}}: 50.
\end{itemize}
All other text is verbatim, with non-ASCII glyphs rendered in ASCII (an em dash as \texttt{-{}-},
$\leq$ as \texttt{<=}, and a left-pointing marker as \texttt{<-}). The returned predicates are then
compiled, screened and assigned to band pools in code (Section~\ref{sec:method:screen}), not by the
LLM.

\begin{lstlisting}[basicstyle=\ttfamily\scriptsize,breaklines=true,columns=fullflexible,keepspaces=true,numbers=none,frame=single,captionpos=b,caption={Rule-proposal prompt template used for every shot (anonymized placeholders, ASCII).},label={lst:prompt}]
{domain primer}

---

{model digest}

---

{SHAP gap}

---

## Task: Adversarial Rule Generation

You are a red-team analyst. The LightGBM model already captures fraud patterns
anchored on its **forbidden features** (high SHAP for fraud, listed above).
Your job is to find fraud patterns the model systematically misses by reasoning
in the **target feature** dimensions it underweights.

Work through the following three steps in your head before writing the output.

### Step 1 -- Adversarial evasion reasoning

Think like a fraudster who knows LGB monitors the forbidden features closely.
Design fraud behaviours that primarily exploit the target features
(SHAP <= 0, listed above) as the main discriminating anchor.
You may also include forbidden features as secondary filters when the combination
reveals a behavioural archetype that the target features alone would not isolate --
for example, a target-feature profile that only becomes high-precision when
narrowed by a forbidden-feature threshold.
For each candidate pattern, ask: *why would a fraudster naturally exhibit this
combination? Does it correspond to a real behavioural archetype?*

### Step 2 -- Data grounding

For each evasion pattern, check the decile/category tables in the SHAP Gap
section. Look for target-feature groups with **Lift > 1.5 (<-)**. Prefer
patterns that are both theoretically coherent AND supported by elevated lift
in the data. If no <- appears for a candidate feature, weigh it lower.

### Step 3 -- Predicate construction

Translate the evasion pattern into a conjunction predicate using:
- Feature names exactly as listed above (no invented names).
- Operators: `>`, `<`, `>=`, `<=`, `==`, `isna`, `notna`.
- Percentile tokens: `p5`, `p10`, `p15`, `p20`, `p25`, `p30`, `p40`,
  `p50`, `p60`, `p70`, `p75`, `p80`, `p85`, `p90`, `p95`, `p99`.
- Join conditions with `AND` only (no top-level OR).

---

### Output format

Produce exactly **{number of rules}** rules as a JSON array. Each item:

```json
[
  {
    "id": "R_001",
    "name": "short label (5 words max)",
    "evasion_reasoning": "why a fraudster shows this pattern in target features; if forbidden features appear as secondary filters, explain why that narrows to a distinct archetype",
    "data_evidence": "specific lift values or decile ranges from the table that support this",
    "target_features_used": ["feat_a", "feat_b"],
    "forbidden_features_used": ["feat_c"],
    "predicate": "feat_a > p80 AND feat_b == 3 AND feat_c > p90"
  }
]
```

Constraints:
- Use ONLY feature names that appear verbatim in the digest or SHAP gap above.
- The primary anchor condition(s) must be target features (SHAP <= 0). Forbidden features may appear as secondary filters in a conjunction.
- Each rule must address a distinct evasion archetype -- no near-duplicates.
- Rank from strongest evidence + clearest evasion logic to weakest.

Respond with the JSON array only -- no explanation, no markdown fences outside the JSON.
\end{lstlisting}

%
\section{Full Pipeline Ablation}
\label{sec:app_pipeabl}

Table~\ref{tab:pipeabl_full} reports both metrics across all eight evaluation splits.
Conditions denote removing components from \textbf{Full} (complete pipeline: all three backbones,
per-band calibration, per-shot logit ensemble): $-$PerBand replaces per-band with global
calibration; $-$Ens uses a merged deduplicated rule pool with no per-shot ensemble;
$-$Ens\,$-$PerBand combines both; $-$LogLift scores each shot by its rule fire rate instead of the
log-lift overlay; Rules-only uses the merged pool's fire rate. Every condition is evaluated like
Full, with its own cut-off table fitted on the adaptation window.

\begin{table*}[t]
\centering
\footnotesize
\caption{Full pipeline ablation: avg$\Delta$R@P (top) and avg$\Delta$P@R (bottom), \%, per split.
  \textbf{Full} = complete pipeline; each other column removes components
  (definitions in Appendix~\ref{sec:app_pipeabl}).}
\label{tab:pipeabl_full}
\begin{tabular}{llrrrrrr}
\toprule
Dataset & Split   & \textbf{Full} & $-$PerBand & $-$Ens & $-$Ens\,$-$PerBand & $-$LogLift & Rules-only \\
\midrule
\multicolumn{8}{@{}l}{\textit{avg$\Delta$R@P}} \\
D-Industrial & OOT & $\mathbf{+2.26\%}$ & $-2.41\%$ & $-2.48\%$ & $-46.39\%$ & $-35.73\%$ & $-24.76\%$ \\
BAF & OOT-1 & $-0.10\%$ & $-0.42\%$ & $-21.42\%$ & $-22.87\%$ & $-26.02\%$ & $-25.76\%$ \\
BAF & OOT-2 & $0.00\%$ & $-0.38\%$ & $-20.44\%$ & $-20.97\%$ & $-25.46\%$ & $-25.16\%$ \\
BAF & Var-I & $\mathbf{+0.04\%}$ & $-0.09\%$ & $-17.45\%$ & $-19.12\%$ & $-21.46\%$ & $-21.04\%$ \\
BAF & Var-II & $\mathbf{+0.27\%}$ & $-0.06\%$ & $-19.96\%$ & $-24.10\%$ & $-32.92\%$ & $-32.10\%$ \\
BAF & Var-III & $\mathbf{+0.04\%}$ & $-0.20\%$ & $-17.85\%$ & $-21.41\%$ & $-26.79\%$ & $-26.40\%$ \\
BAF & Var-IV & $\mathbf{+0.04\%}$ & $-0.20\%$ & $-17.85\%$ & $-21.41\%$ & $-26.79\%$ & $-26.40\%$ \\
BAF & Var-V & $-0.13\%$ & $-0.70\%$ & $-21.97\%$ & $-24.63\%$ & $-29.18\%$ & $-28.44\%$ \\
\midrule
\multicolumn{8}{@{}l}{\textit{avg$\Delta$P@R}} \\
D-Industrial & OOT & $\mathbf{+1.60\%}$ & $-2.99\%$ & $-5.11\%$ & $-41.54\%$ & $-31.45\%$ & $-24.11\%$ \\
BAF & OOT-1 & $\mathbf{+0.07\%}$ & $-0.25\%$ & $-15.46\%$ & $-19.38\%$ & $-26.84\%$ & $-26.00\%$ \\
BAF & OOT-2 & $\mathbf{+0.06\%}$ & $-0.21\%$ & $-16.18\%$ & $-19.68\%$ & $-28.15\%$ & $-27.15\%$ \\
BAF & Var-I & $\mathbf{+0.08\%}$ & $+0.25\%$ & $-15.73\%$ & $-17.79\%$ & $-26.75\%$ & $-26.37\%$ \\
BAF & Var-II & $\mathbf{+0.30\%}$ & $+0.10\%$ & $-13.45\%$ & $-16.36\%$ & $-23.78\%$ & $-22.71\%$ \\
BAF & Var-III & $\mathbf{+0.28\%}$ & $+0.26\%$ & $-9.25\%$ & $-11.49\%$ & $-16.42\%$ & $-15.69\%$ \\
BAF & Var-IV & $\mathbf{+0.28\%}$ & $+0.26\%$ & $-9.25\%$ & $-11.49\%$ & $-16.42\%$ & $-15.69\%$ \\
BAF & Var-V & $\mathbf{+0.02\%}$ & $+0.27\%$ & $-11.20\%$ & $-12.10\%$ & $-17.61\%$ & $-16.74\%$ \\
\bottomrule
\end{tabular}
\end{table*}

%
\section{Full LLM-Source Ablation}
\label{sec:app_srcabl}

Table~\ref{tab:srcabl_full} reports both metrics across all eight evaluation splits for all seven
combinations of the three LLM backbones (full pipeline). No combination dominates uniformly:
LLM-MoE alone gives the highest avg$\Delta$R@P on D-Industrial ($+3.88\%$) and
LLM-SlW$+$LLM-MoE the highest of the pairs ($+2.80\%$), but both are negative on most BAF
splits, while LLM-CoT---slightly negative on D-Industrial---is the best single source on every BAF
split. All~3 remains strongly positive on D-Industrial and avoids the worst corners on BAF.

\begin{table*}[t]
\centering
\footnotesize
\setlength{\tabcolsep}{4pt}
\caption{Full LLM-source ablation: avg$\Delta$R@P (top) and avg$\Delta$P@R (bottom), \%, per split,
  all seven backbone combinations (full pipeline). \textbf{All~3} = our default ensemble.
  CoT/SlW/MoE abbreviate LLM-CoT/LLM-SlW/LLM-MoE (Section~\ref{sec:setup}).}
\label{tab:srcabl_full}
\begin{tabular}{llrrrrrrr}
\toprule
Dataset & Split & CoT & SlW & MoE & CoT+SlW & CoT+MoE & SlW+MoE & \textbf{All~3} \\
\midrule
\multicolumn{9}{@{}l}{\textit{avg$\Delta$R@P}} \\
D-Industrial & OOT & $-0.28\%$ & $+2.27\%$ & $+3.88\%$ & $+1.14\%$ & $+2.36\%$ & $+2.80\%$ & $\mathbf{+2.26\%}$ \\
BAF & OOT-1 & $-0.12\%$ & $-0.21\%$ & $-0.18\%$ & $+0.06\%$ & $-0.03\%$ & $-0.33\%$ & $-0.10\%$ \\
BAF & OOT-2 & $-0.08\%$ & $-0.41\%$ & $-0.26\%$ & $+0.14\%$ & $+0.08\%$ & $-0.29\%$ & $0.00\%$ \\
BAF & Var-I & $+0.21\%$ & $-0.07\%$ & $-0.22\%$ & $+0.27\%$ & $+0.16\%$ & $-0.20\%$ & $\mathbf{+0.04\%}$ \\
BAF & Var-II & $+0.37\%$ & $+0.20\%$ & $+0.19\%$ & $+0.27\%$ & $+0.33\%$ & $+0.18\%$ & $\mathbf{+0.27\%}$ \\
BAF & Var-III & $+0.23\%$ & $-0.22\%$ & $-0.07\%$ & $+0.27\%$ & $+0.26\%$ & $-0.18\%$ & $\mathbf{+0.04\%}$ \\
BAF & Var-IV & $+0.23\%$ & $-0.22\%$ & $-0.07\%$ & $+0.27\%$ & $+0.26\%$ & $-0.18\%$ & $\mathbf{+0.04\%}$ \\
BAF & Var-V & $+0.10\%$ & $-0.51\%$ & $-0.36\%$ & $+0.02\%$ & $-0.12\%$ & $-0.42\%$ & $-0.13\%$ \\
\midrule
\multicolumn{9}{@{}l}{\textit{avg$\Delta$P@R}} \\
D-Industrial & OOT & $-0.15\%$ & $+1.56\%$ & $+3.09\%$ & $+0.78\%$ & $+1.80\%$ & $+2.34\%$ & $\mathbf{+1.60\%}$ \\
BAF & OOT-1 & $+0.12\%$ & $+0.06\%$ & $-0.09\%$ & $+0.27\%$ & $+0.06\%$ & $-0.05\%$ & $\mathbf{+0.07\%}$ \\
BAF & OOT-2 & $+0.09\%$ & $-0.17\%$ & $-0.14\%$ & $+0.24\%$ & $+0.06\%$ & $-0.04\%$ & $\mathbf{+0.06\%}$ \\
BAF & Var-I & $+0.12\%$ & $+0.17\%$ & $+0.04\%$ & $+0.19\%$ & $+0.13\%$ & $+0.44\%$ & $\mathbf{+0.08\%}$ \\
BAF & Var-II & $+0.25\%$ & $+0.13\%$ & $+0.13\%$ & $+0.23\%$ & $+0.16\%$ & $+0.11\%$ & $\mathbf{+0.30\%}$ \\
BAF & Var-III & $+0.26\%$ & $+0.18\%$ & $+0.41\%$ & $+0.25\%$ & $+0.33\%$ & $+0.37\%$ & $\mathbf{+0.28\%}$ \\
BAF & Var-IV & $+0.26\%$ & $+0.18\%$ & $+0.41\%$ & $+0.25\%$ & $+0.33\%$ & $+0.37\%$ & $\mathbf{+0.28\%}$ \\
BAF & Var-V & $+0.01\%$ & $-0.13\%$ & $-0.28\%$ & $+0.09\%$ & $+0.09\%$ & $+0.01\%$ & $\mathbf{+0.02\%}$ \\
\bottomrule
\end{tabular}
\end{table*}

%
\section{All Operating Points}
\label{sec:app_perband}

Table~\ref{tab:perband_all} lists, for every cut-off $t_b$ and split, production ($s\ge t_b$) and
the all-3 ensemble overlay ($\tilde s_b\ge t_b$) side by side on OOT data---precision and recall---with
the $\Delta$R@P and $\Delta$P@R checks at that cut-off. The two need not alert the same volume on
OOT; the checks compare them at equal precision and equal recall instead. Var-IV shares Var-III's evaluation population.

\begin{table*}[t]
\centering
\scriptsize
\setlength{\tabcolsep}{4pt}
\caption{All operating points, all-3 ensemble: production $s\ge t_b$ vs.\ overlay
  $\tilde s_b\ge t_b$ on OOT data (precision P\%, recall R\%), and $\Delta$R@P / $\Delta$P@R
  (\%) at each cut-off. Every parameter is fixed on the adaptation window.}
\label{tab:perband_all}
\begin{tabular}{llrrrrrr}
\toprule
& & \multicolumn{2}{c}{Production} & \multicolumn{2}{c}{Overlay} & & \\
\cmidrule(lr){3-4}\cmidrule(lr){5-6}
Split & $t_b$ & P\% & R\% & P\% & R\% & $\Delta$R@P & $\Delta$P@R \\
\midrule
D-Ind. OOT & 998 & 56.48 & 23.47 & 57.51 & 24.37 & $+2.57\%$ & $+1.03\%$ \\
 & 997 & 49.07 & 35.67 & 50.30 & 38.00 & $+3.31\%$ & $+2.04\%$ \\
 & 996 & 43.29 & 45.80 & 44.06 & 47.67 & $+1.88\%$ & $+1.66\%$ \\
 & 995 & 37.99 & 52.41 & 38.73 & 54.00 & $+2.41\%$ & $+1.91\%$ \\
 & 994 & 33.99 & 58.12 & 34.58 & 59.27 & $+1.55\%$ & $+1.96\%$ \\
 & 993 & 30.32 & 62.86 & 30.91 & 64.00 & $+1.84\%$ & $+1.02\%$ \\
\midrule
BAF OOT-1 & 999 & 53.15 & 4.12 & 53.95 & 4.05 & $+0.21\%$ & $+0.30\%$ \\
 & 995 & 34.89 & 13.88 & 34.68 & 13.53 & $-0.59\%$ & $-0.11\%$ \\
 & 990 & 28.15 & 22.40 & 28.55 & 21.74 & $-0.31\%$ & $-0.34\%$ \\
 & 985 & 23.76 & 28.35 & 24.39 & 28.00 & $+0.94\%$ & $+0.52\%$ \\
 & 980 & 20.77 & 33.07 & 21.08 & 33.31 & $+0.94\%$ & $+0.60\%$ \\
 & 975 & 18.94 & 37.68 & 19.00 & 37.05 & $-1.68\%$ & $-0.45\%$ \\
 & 970 & 17.35 & 41.42 & 17.56 & 40.96 & $-0.21\%$ & $-0.01\%$ \\
\midrule
BAF OOT-2 & 999 & 52.83 & 3.89 & 52.68 & 3.75 & $+0.31\%$ & $+0.71\%$ \\
 & 995 & 37.20 & 13.83 & 36.60 & 13.24 & $-1.53\%$ & $-0.47\%$ \\
 & 990 & 29.74 & 21.86 & 30.08 & 21.09 & $-0.35\%$ & $-0.22\%$ \\
 & 985 & 25.23 & 27.52 & 25.54 & 26.93 & $+1.08\%$ & $+0.36\%$ \\
 & 980 & 22.57 & 32.97 & 22.64 & 32.80 & $+0.69\%$ & $+0.10\%$ \\
 & 975 & 20.51 & 37.32 & 20.70 & 36.80 & $-0.14\%$ & $-0.01\%$ \\
 & 970 & 19.07 & 41.17 & 19.30 & 40.93 & $-0.10\%$ & $-0.08\%$ \\
\midrule
BAF Var-I & 999 & 56.25 & 3.13 & 57.79 & 3.09 & $-0.03\%$ & $-0.25\%$ \\
 & 995 & 33.98 & 10.88 & 34.50 & 10.63 & $+0.07\%$ & $+0.07\%$ \\
 & 990 & 28.05 & 17.72 & 28.91 & 17.48 & $0.00\%$ & $+0.56\%$ \\
 & 985 & 24.64 & 23.45 & 25.22 & 23.11 & $+0.59\%$ & $+0.39\%$ \\
 & 980 & 21.85 & 28.11 & 22.23 & 28.04 & $+0.56\%$ & $+0.30\%$ \\
 & 975 & 19.65 & 32.00 & 19.97 & 31.62 & $-0.28\%$ & $-0.17\%$ \\
 & 970 & 17.99 & 35.27 & 18.03 & 34.64 & $-0.63\%$ & $-0.33\%$ \\
\midrule
BAF Var-II & 999 & 51.84 & 6.36 & 53.04 & 6.36 & $-0.21\%$ & $+1.18\%$ \\
 & 995 & 31.96 & 20.29 & 32.32 & 20.01 & $-0.28\%$ & $-0.41\%$ \\
 & 990 & 24.67 & 30.09 & 25.27 & 29.78 & $+0.21\%$ & $+0.36\%$ \\
 & 985 & 20.40 & 36.69 & 20.91 & 36.28 & $+1.04\%$ & $+0.35\%$ \\
 & 980 & 18.03 & 42.49 & 18.24 & 42.43 & $+0.03\%$ & $+0.13\%$ \\
 & 975 & 16.03 & 46.39 & 16.39 & 46.53 & $+1.04\%$ & $+0.45\%$ \\
 & 970 & 14.58 & 50.14 & 14.76 & 49.69 & $+0.07\%$ & $+0.05\%$ \\
\midrule
BAF Var-III & 999 & 36.59 & 4.52 & 38.27 & 4.68 & $+0.04\%$ & $+1.59\%$ \\
 & 995 & 21.17 & 15.12 & 21.60 & 15.05 & $0.00\%$ & $+0.01\%$ \\
 & 990 & 16.54 & 23.87 & 16.96 & 23.57 & $-0.72\%$ & $+0.02\%$ \\
 & 985 & 13.95 & 30.13 & 14.32 & 29.83 & $+1.09\%$ & $+0.26\%$ \\
 & 980 & 12.45 & 35.82 & 12.53 & 35.52 & $-0.15\%$ & $-0.08\%$ \\
 & 975 & 11.10 & 39.63 & 11.30 & 39.59 & $+0.08\%$ & $+0.19\%$ \\
 & 970 & 10.11 & 43.06 & 10.22 & 42.61 & $-0.08\%$ & $-0.01\%$ \\
\midrule
BAF Var-IV & 999 & 36.59 & 4.52 & 38.27 & 4.68 & $+0.04\%$ & $+1.59\%$ \\
 & 995 & 21.17 & 15.12 & 21.60 & 15.05 & $0.00\%$ & $+0.01\%$ \\
 & 990 & 16.54 & 23.87 & 16.96 & 23.57 & $-0.72\%$ & $+0.02\%$ \\
 & 985 & 13.95 & 30.13 & 14.32 & 29.83 & $+1.09\%$ & $+0.26\%$ \\
 & 980 & 12.45 & 35.82 & 12.53 & 35.52 & $-0.15\%$ & $-0.08\%$ \\
 & 975 & 11.10 & 39.63 & 11.30 & 39.59 & $+0.08\%$ & $+0.19\%$ \\
 & 970 & 10.11 & 43.06 & 10.22 & 42.61 & $-0.08\%$ & $-0.01\%$ \\
\midrule
BAF Var-V & 999 & 40.36 & 5.05 & 40.54 & 5.09 & $+0.04\%$ & $+0.18\%$ \\
 & 995 & 22.43 & 16.48 & 22.76 & 16.29 & $-0.04\%$ & $-0.07\%$ \\
 & 990 & 17.52 & 25.45 & 17.76 & 25.00 & $+0.41\%$ & $+0.11\%$ \\
 & 985 & 14.74 & 31.86 & 15.03 & 31.18 & $+0.04\%$ & $+0.01\%$ \\
 & 980 & 13.16 & 37.78 & 13.13 & 37.18 & $-1.09\%$ & $-0.07\%$ \\
 & 975 & 11.82 & 42.16 & 11.98 & 41.67 & $+0.34\%$ & $+0.15\%$ \\
 & 970 & 10.79 & 45.59 & 10.89 & 45.06 & $-0.64\%$ & $-0.17\%$ \\
\bottomrule
\end{tabular}
\end{table*}

%
\section{D-Industrial per OOT Month}
\label{sec:app_month}

Table~\ref{tab:month} breaks the D-Industrial results of the all-3 ensemble down by OOT month, with
the cut-offs and every table fixed on the adaptation window and precision, recall, and the checks
computed within each month. March covers only its first seven days, so its per-cut-off values rest on
few frauds.

\begin{table*}[t]
\centering
\footnotesize
\caption{D-Industrial per OOT month, all-3 ensemble: $\Delta$R@P (top) and $\Delta$P@R (bottom) at each
  production cut-off $t_b$ and their average (\%).}
\label{tab:month}
\begin{tabular}{lrrrrrrr}
\toprule
Month & 998 & 997 & 996 & 995 & 994 & 993 & Avg \\
\midrule
\multicolumn{8}{@{}l}{\textit{$\Delta$R@P}} \\
Jan 2026 & $+6.10\%$ & $+1.72\%$ & $+0.47\%$ & $+2.27\%$ & $+1.56\%$ & $+2.27\%$ & $+2.40\%$ \\
Feb 2026 & $+2.29\%$ & $+5.01\%$ & $+3.48\%$ & $+3.16\%$ & $+0.98\%$ & $+1.74\%$ & $+2.77\%$ \\
Mar 2026 & $-0.79\%$ & $+4.74\%$ & $+3.56\%$ & $+1.98\%$ & $+0.79\%$ & $-0.40\%$ & $+1.65\%$ \\
\midrule
\multicolumn{8}{@{}l}{\textit{$\Delta$P@R}} \\
Jan 2026 & $+2.72\%$ & $+2.37\%$ & $+0.88\%$ & $+1.64\%$ & $+2.35\%$ & $+1.10\%$ & $+1.84\%$ \\
Feb 2026 & $+1.13\%$ & $+2.96\%$ & $+2.78\%$ & $+2.08\%$ & $+0.84\%$ & $+0.84\%$ & $+1.77\%$ \\
Mar 2026 & $-0.28\%$ & $+2.00\%$ & $+1.91\%$ & $+1.64\%$ & $+1.78\%$ & $-0.61\%$ & $+1.07\%$ \\
\bottomrule
\end{tabular}
\end{table*}


\end{document}
